\section{Analysis}
\label{sec:analysis}

\subsection{Data and Variables}
\label{subsec:data}

\subsubsection{Data Source}

In the empirical analysis that follows, I consider data obtained from the Arts and Culture Supplement to the National Social Survey fielded in 2002 ($N=1,002$) and 2004 ($N=801$) \citep{NSS2002, NSS2004}.  The NSS is a weighted (by age, education, marital status, gender, race, ethnicity) probability sample of respondents residing the United States collected through computer-aided telephone interviews.   The NSS is particularly appropriate for my purposes because both waves contain information not only about culture consumption experiences, but also about the likelihood that individuals engage in informal social interaction featuring cultural content taken from consumption in the same domains.  Because the items of interest were fielded in both waves, I pool individuals for which I have complete data on all socio-demographic variables from the 2002 ($N=1002$) and 2004 ($N=801$) surveys, resulting in a total combined sample of 1,803 respondents.

\subsubsection{Key response and control variables}
As part of the survey, respondents were asked to report (during the past twelve months), ``approximately how often have you had a discussion with someone you know'' about (1) a movie, (2) a musical performance, live or recorded, (3) a book, a poem or a story, (4) an art show or a work of art, and (5) a play or other dramatic performance.  These are the main empirical indicators of culture talk and thus will serve as the response variables in the analyses that follow.  Respondents were also asked to report how often (in the last year) they had been to the movies, listened to recorded music at home, read a book, had gone to an art show or a museum, or had gone to any live performances.  Because I am interested in accounting for residual socio-demographic differences in culture talk net of consumption, these serve as the main control variables in the analysis.  For both sets of responses, respondents were asked to choose from the following response categories: 1) very often, 2) fairly often, 3) not very often, 4) or not at all.\footnote{How does the NSS compare to other Arts participation Surveys?  When it comes to book reading, the 2002 NSS estimates that about 55.7\% of respondents report reading books ``Fairly Often'' or `` Very Often.''  The corresponding estimate for the 2002 SPPA (Survey for Public Participation in the Arts) is about 58\%.  When it comes to going to an art museum, about 12\% of NSS respondents report going ``Fairly Often'' and about 3.7\% of respondents report going ``Very Often''; these correspond well with SPPA 2002 data based on more specific questions on attendance frequencies:  13.4\% of SPPA respondents report having gone to the museum between 2 and 5 times in the last year, and 3.2\% report having gone more than five times.  In the NSS, about 9\% of respondents report going to a live performance ``Very Often.''  In the 2002 SPPA about 7.3\% of respondents report having gone out to a Jazz, Classical, Opera, Musical, Ballet or Dance performance between 2 and 5 times in the last year.  In all, the NSS appears to do a good job of capturing rates of arts participation that are comparable to that obtained from other National Surveys.}
 
\subsection{Analytic Strategy}
\subsubsection{Propensity of culture talk}
As noted above, one obvious answer to the question of the socio-demographic structuration of culture talk suggests that this outcome should be most strongly predicted by culture-consumption. Insofar as it is direct experiences with cultural goods and performances that come to serve as the source-content for culture talk episodes we should expect that individuals who are active culture consumers should also evince a higher propensity to engage in culture talk. However, my main interest is status-linked differences in culture talk propensity that are not directly traceable to diversity in consumption behavior.  Thus, I specify an ordered probit model \citep{long_freese06} that holds constant self-reported consumption frequency (as given by the above indicators), in order to ascertain whether there is residual variation across status groups (defined by levels of education) in culture talk propensity \emph{within-levels} of self-reported consumption.  This model takes the form:

\begin{equation}
E(P[Y_{it} \leq k])= \alpha_k + \gamma_t E_i + \delta_t C_{it} + \sum_k\beta_{kt}X_{ik} 
\label{eq:eq1}
\end{equation}•

Where $Y_{it}$ is the frequency of culture talk for individual $i$  in the $t^{th}$ cultural domain, $E_{i}$ is the educational attainment of that person, $C_{it}$ is the self-reported consumption frequency for each individual in that domain.  The effect of these last two  variables expected to be large and positive (and $\gamma_t>0$ and $\delta_t>0$). Finally, $\beta_k$ is an $i \times k$ vector of coefficient estimates corresponding to the effect of the other predictors included in the model ($X_{ik}$); this is the \emph{additive model}.  If there is residual status-linked variation in culture talk net of consumption then we should expect that $\gamma \neq 0$, after holding constant $W_{it}$.  

Alternatively, we may expect that the status effect should itself vary according to the level of self-reported consumption.  This would be the case, for instance, if persons of high-education engage in culture talk at higher rates than persons of low-education even when they consume culture at the same rates. I examine the validity of this hypothesis by specifying a model of the form:

\begin{equation}
E(P[Y_{it} \leq k])= \alpha_k + \gamma_t E_i + \sum_{j=2}^{4}  \delta_{jt} C_{ijt} + \sum_{j=2}^{4} \rho_{jt} C_{ijt}E_{ij} + \sum_k\beta_{kt}X_{ik}
\label{eq:eq2}
\end{equation}•

This is the \emph{multiplicative} model. In this model I enter the frequency of consumption as (4-1=3) binary indicators ($C_{ijt}$) and interact each with the education variable. As \citet{ai_norton03}, have recently noted, the empirical validity of hypothesis 2---which says that the effect of education varies across self-reported consumption status---cannot be verified simply by looking at the \emph{t}-statistics corresponding to each of the $\rho_{jt}$ effects. This is due to the fact that in non-linear (e.g. logit or probit) models the actual interaction effects depend on the value of all other variables in the model any may be present even if $\rho_{jt}=0$. Instead, after estimating the model, I compute the marginal effect of moving from the lowest to the highest education level on the probability of engaging in culture talk ``very often.''  I then graph the corresponding point estimate and 90\% confidence interval across groups of respondents defined by self-reported frequency levels.  The presence of a statistically discernible upward trend in the marginal effects estimates would constitute positive evidence of the expected interaction between status and frequency on consumption for that domain. 

\subsubsection{Indicators of socio-economic status}
I estimate five ordered probit models (one for each cultural domain) of the additive and multiplicative form with the culture talk indicator as the response variable, a series of socio-demographic variables as the predictor variables and the frequency of consumption for that cultural domain as the main control variable. The predictors were chosen  using a combination of criteria geared towards substantive relevance and predictive power. I was also limited by the availability of socio-demographic indicators in the survey: for instance, the NSS contains no information on respondent's current occupation or family background (both of which are theoretically relevant dimensions of social differentiation). 

The main indicator of socio-economic status, as indicated above, is educational attainment and, secondarily, family income.  The education variable is coded as five ordered categories (from high-school dropout to graduate or professional degree) and the income variable is coded in six ordered categories (from less than 20K a year to 100K or more).\footnote{Family income information was not available for about one fifth of the cases (20.8\%) across the two survey waves.  For these cases ($N=375$) I imputed their income using the multiple imputation procedure in Stata 11.  Essentially this involves using an ordered probit model to predict the missing income values using all remaining socio-demographic information for that respondent (e.g. education, age, marital status and selected interactions.} 

\subsubsection{Additional socio-demographic predictors}
The models presented below also include the following vector of socio-demographic covariates:  (1) respondent's age (in seven ordered categories),  (2) respondent's gender, and (3) urban Location (urban versus other location).  I experimented with models that included a vector of dummy variables coding for race/ethnic identification and marital status.  These variables had for the most part non-statistically significant effects in all of the equations.  To avoid over-parameterizing the models and for the sake of parsimony, I excluded these variables from the presentation of the results. Tables showing the coefficient estimates and t-statistics for these variables are available upon request from the author.  Descriptive statistics corresponding to each of the predictor variables included in the models are presented in Table~\ref{tab:desc} of the Appendix.  

For the sake of efficiency, simplicity, and to achieve comparability in the magnitude of the effects across the different types of predictors in the presentation of the results, I enter all of the ordered categorical predictors as numeric inputs into the various equations (summarizing the effect of education, for instance, in a single coefficient rather than ($K-1$ indicator variables, where $K$ is the number of ordered categories).  This linearity constraint results in some loss of information, but the substantive story remains the same regardless of the strategy I pursue (results available on request).  To facilitate comparison of the magnitude of the effects across inputs, I standardized each of the predictors by subtracting the mean and dividing by two times the standard deviation as recommended by \citet{gelman08}.  This makes the interpretation of  the (ordered categorical) numeric predictors such as education, income, and age straightforward: the coefficient estimate indicates the effect of moving from the lowest to the highest value, while the binary predictors (gender, urban and marital status) retain their usual interpretation.


